\section{Background}
\label{sec:background}

\subsection{Disaggregated memory and paging}
In a disaggregated deployment, the compute node holds only a portion of the memory that its applications actually address---perhaps one fifth of the total working set---while the remainder lives on separate memory nodes. Three broad design families have been explored. Library- or runtime-level approaches, such as AIFM and Semeru, surface remote memory through new programming abstractions and therefore need application modifications \cite{aifm,semeru}. Hardware proposals extend the memory bus itself, as the Compute Express Link standard does \cite{cxl}. Paging-based designs take a third path: they reuse the existing virtual-memory machinery by marking non-resident pages as non-present, letting the page-fault handler perform the remote fetch transparently \cite{infiniswap,fastswap,leap,hermit,dilos}. This last family requires no changes to applications, which is why it is the focus of the paper we examine.

A fault in this setting always follows the same sequence of steps. The handler allocates a local frame (evicting another page if necessary), posts an RDMA read for the missing 4~KiB page, waits for the read to complete, maps the page into the process's address space, and returns to the application. Memory pressure is managed by a background reclaim path that writes dirty pages back and recycles their frames. Of all these steps, only the waiting is relevant to the yield-versus-spin question, so the rest of this report focuses on it exclusively.

\subsection{RDMA fetch latency and its components}
A one-sided RDMA read bypasses the remote CPU entirely; its latency breaks down into work at the local NIC, wire propagation, and work at the remote NIC. Kalia et al.\ give a detailed account of how queue pairs (QPs), doorbells, and completion queues shape the throughput available to a host \cite{kalia}. For our purposes, three properties are what matter. First, a single fetch completes within a few microseconds when the fabric is unloaded. Second, commands on a single QP are processed in order, so bursts of requests on one QP pile up behind each other. Third, all QPs compete for the same physical link and NIC pipeline, capping the aggregate request rate. Our simulator captures exactly these three properties and nothing beyond them (Section~\ref{sec:model}).

\subsection{Two ways to wait}
When a worker's request faults, it faces a choice between two strategies.

\emph{Busy-waiting} keeps the faulting request on the core throughout the fetch. The worker polls the completion queue until the read is done. No scheduler involvement is needed and there are no context switches, so the added overhead is virtually zero. The drawback is that the core is tied up and unavailable to serve other requests for the full duration of the fetch.

\emph{Yielding} takes the opposite approach: the request's context is saved, control is returned to the scheduler, and another request runs in the meantime. Once the completion arrives, the scheduler resumes the original request. The core stays productive during the fetch, but every page fault now carries the cost of two context switches plus the surrounding scheduling work. In a standard OS kernel, those costs dwarf the few-microsecond fetch, making yielding a net loss. The core idea behind Adios is that user-level threads running inside a single address space can shrink the switch cost enough that yielding becomes profitable again \cite{adios}. Switching stacks without crossing a kernel boundary---as systems like Capriccio, Arachne, and Shenango do---keeps latency at the microsecond level \cite{capriccio,arachne,shenango}.

\subsection{Why spinning hurts the tail}
Imagine $N$ workers pulling from a shared FIFO queue. When every fault pins a worker to the core for the entire fetch time $t_f$, the mean per-request occupancy is large, and the system's hard throughput ceiling is $N$ divided by that mean. Below this ceiling the queue stays short. As the offered load closes in on the ceiling, an arriving request must queue behind workers that are stuck waiting for RDMA---this is head-of-line (HOL) blocking, the same phenomenon found in packet switches and microsecond-scale schedulers \cite{zygos,shinjuku}: short local-hit requests get stuck behind long remote-miss requests, and queueing delay rapidly dominates the tail. Preemptive schedulers can slice a long request into small pieces \cite{shinjuku,concord}, but a waiting worker gains nothing from preemption because the wait itself is not computation. The Adios authors verified this directly: adding a 5~$\mu$s preemption interval to a spinning system made no difference to the tail \cite{adios}.

\subsection{Why spinning wastes bandwidth}
A spinning worker can keep at most one RDMA fetch in flight. With eight workers, that is an absolute ceiling of eight concurrent fetches, regardless of how much more the NIC and link could handle. A yielding worker, by contrast, may have many fetches outstanding simultaneously. The practical result is that the busy-waiting system exhausts its \emph{cores} while the network link sits half idle, whereas the yielding system keeps the fetch queue fuller and shifts the bottleneck to the link instead. Section~\ref{sec:bounds} puts precise numbers behind this argument.

\subsection{Unikernels and lightweight threads}
A unikernel bundles the application and OS services into a single image that runs in one address space under one protection domain \cite{unikernels}. OSv, Unikraft, and Unikernel Linux are well-known examples \cite{osv,unikraft,ukl}. Without a user/kernel boundary, a page-fault handler can hand control to the scheduler through an ordinary function call rather than a system call. Adios exploits this property to push the yield cost down to tens of cycles \cite{adios}. In our simulator we do not model the unikernel at all; instead, we expose the switch cost as a free parameter $t_s$ so that we can study how the scheduling decision scales with it.
